\documentclass{article}
\usepackage{amsmath, amssymb}

\title{Analysis - I}
\author{Virat Dixit}
\date{Day 1}

\begin{document}

\maketitle

\section{What is Analysis?}

Analysis is the rigorous study of continuous change, limits, and the behaviour of functions, sequences and series, dealing with real and complex numbers, so as to understand highly abstract fields like calculus more deeply.

\subsection*{Key Subfields}
\begin{enumerate}
    \item Real Analysis
    \item Complex Analysis
    \item Functional Analysis
\end{enumerate}

In these notes we will explore Real Analysis following Tao's approach.

\section{Real Analysis}

Real Analysis is the theoretical foundation which underlies calculus, which is the collection of computational algorithms which one uses to manipulate functions.

\medskip
Broadly, Real Analysis is the study of real numbers and their topological properties, sequences and series of real numbers, and properties of real valued functions.

\section{Numbers}

The natural numbers $\{0, 1, 2, 3, \ldots\}$ are the most primitive of the number systems, but they are used to build the integers, which in turn are used to build the rationals, and further the rationals are used to build the real numbers, which are used to build the complex numbers.

\subsection*{The Peano Axioms}

Peano axioms were first laid out by Giuseppe Peano (1858-1932). This is the standard way to define Natural Numbers.

\section{Defining Natural Numbers}

\textbf{Definition 1 (Informal)} - A natural number is any element of the set
\[
    \mathbb{N} := \{0, 1, 2, 3, 4 \ldots \}
\]
which is the set of all numbers created by starting with 0 and then counting forward indefinitely. We call $\mathbb{N}$ the set of natural numbers.

\subsection*{Analysing the Definition}

\begin{itemize}
    \item In this text, the natural numbers includes 0 (which makes it whole numbers), but as the text mentions, some mathematical conventions start natural numbers at 1 instead of 0.

    \item The text recognizes $\mathbb{N}$ as $\{0, 1, 2, \ldots\}$ and positive integers $\mathbb{Z}^+$ as $\{1, 2, 3, 4 \ldots\}$.

    \item Notation: The ``$:=$'' symbol means ``is defined as''.

    \item Terms like ``counting forward indefinitely'' are not mathematically precise. It doesn't explain how you get from one number to another or ensure that you don't loop back to previous numbers.
\end{itemize}

\end{document}
